\section{PA194R: a reassociated 194-state decoder}
\label{sec:pa194r}

PA194R combines the final-emitter decrement proved above with a different
binary association of the axiom-emitter program. It has 194 working
states and 388 transition rows. The complete primitive instruction
graph is identical to PA194F, but 1,111 of its 3,646 source addresses
change. These addresses are the binary paths physically processed by
the decoder, so the construction changes the encoded program and its
shared decoder graph. It is more than a relabeling of PA194F's controls.
The known-zero closing principle and symmetric cached controller remain
the explicitly attributed inherited components.

Of the 67 axiom-reader subtree types, 23 are shared with PA194F and
44 are replaced by different types. The maximum complete source-address
length remains 37. Thus the claim of a different decoder is supported
both by its changed sharing structure and by the separate control-graph
isomorphism check, not merely by a different printed numbering.

The new association retains 67 distinct axiom-reader states before
folding. Together with the 135 fixed controls, the unfolded machine has
202 working states. Seven contextual folds in the axiom reader and the
retained symmetric-controller fold give
\[
                         135+67-7-1=194.
\]
The physical state savings therefore survive the changed program
association. The construction and fold records are reproducible from
\texttt{construction/build\_reassociated.py}; the final machine and
its source paths are retained in \texttt{machines/PA194R/}.

The dense table's SHA-256 is
\begin{quote}\small\ttfamily
\nolinkurl{8a9a8a30773f73b638a7c51da1afca9c9921371aea819d6a21f02d742f673a4e},
\end{quote}
and the native table's SHA-256 is
\begin{quote}\small\ttfamily
\nolinkurl{24539448def6e974d636052dbcb56eeb61d8e06ad0e9924d734c3bd4f0e01aa8}.
\end{quote}

\subsection{Exact source equality and independent literal verification}

Preserving a flattened instruction word would not be sufficient here:
breaks and continues depend on the surrounding binary program structure.
The independent verifier instead parses the delivered source paths and
resolves every successor twice, using the arithmetic-path and interval
algorithms. The results agree at all 3,646 sites, and agree exactly with
PA194F's resolved graph, including both targets of every decrement.
The primitive operation and register operand also agree at every site.
Thus the change is in physical source addresses and decoding; it makes
no change to the already-proved primitive source machine $P_F$.

The standalone verifier
\texttt{verification/check\_reassociated.py} reads the delivered native
transition table directly. It does not import the construction compiler
or use an arithmetic-register oracle. It rederives the four register-list
lemmas from the literal rows and checks all 4,746 primitive outcomes,
including the positive branch at the final emitter decrement. Each
contract quantifies over arbitrary natural register values, an arbitrary
source gap, and an arbitrary finite nonempty list of arbitrary natural
spare values. The contracts certify complete frame transformations and
positive physical durations, with no intermediate halt in a nonhalting
instruction segment.

A separate simulator stores the finite set of tape cells containing one,
so the exterior is exactly the infinite zero tape. From native start
control 74 it reaches the complete clean source frame after 57,768
transitions: native control 84, head $-2$, source site 25, logical
registers $(2,0,0,0,0,0,0,0,0)$, and the recorded 26 spare registers.
The entire finite tape support is compared with the frame encoding,
not just the decoded logical registers. The source path at site 25 is
unchanged by reassociation. The dense table is checked to be the exact
breadth-first renaming of the native table, with dense start 0 and
halting target 194.

Receipts, all symbolic contracts, the rederived list lemmas, the complete
blank seed and the list of changed source addresses are retained in
\texttt{verification/recorded/PA194R/}. The verifier supports
\texttt{--output} for replay into a fresh directory and rejects Python's
optimized mode, which would disable its assertions.

\begin{theorem}[Consistency calibration of PA194R]
For the delivered binary table with its specified blank start,
\[
 \mathrm{PA}\vdash
 H_{\mathrm{PA194R}}
 \ \longleftrightarrow\ H_{P_F}
 \ \longleftrightarrow\
 \exists p\,\operatorname{Prf}_{\mathrm{PA}}
                         (p,\ulcorner0=1\urcorner).
\]
Consequently
\[
 \mathrm{PA}\vdash
 \neg H_{\mathrm{PA194R}}\ \longleftrightarrow\
                       \operatorname{Con}(\mathrm{PA}).
\]
\end{theorem}
\begin{proof}
Exact graph equality identifies PA194R's source machine with $P_F$.
Its independently checked startup supplies the same logical seed.
Concatenating the total positive instruction contracts maps a finite
source halt to a finite literal halt. Conversely, a literal halt at
time $T$ cannot coexist with $T+1$ nonhalting source instructions:
their positive physical segments, preceded by startup, would extend a
deterministic nonhalting literal execution beyond time $T$. This is the
same bounded converse argument established for PA194F. The source
equivalence and inherited syntactic PA calibration then give the displayed
biconditionals. All finite checks and run inductions have the ordinary
mathematical PA-provability scope explained in the preceding sections.
\end{proof}

The changed association need not preserve physical running times or
spare-list lengths. Each complete-frame contract accounts for its own
effects; source equality is what transfers the logical calibration.
As with PA194F, the new result is supported by ordinary proofs and
complete symbolic verification, with the inherited calibration and
symbolic-rule soundness dependencies stated explicitly. No new
end-to-end proof-assistant kernel theorem is asserted.
